% =====================================================
% 题型：等差数列实际应用与分组选择题 (q_seq_pagoda.tex)
% 知识板块：数列、组合配对
% 主思维：结构归纳与套用
% 副思维：特殊值试探、化归转化
% 错因标签：未能利用奇偶性排除、分组配对困难
% 讲义用途：数列结构题
% =====================================================
% 难度：★★ (中档)
% =====================================================
\newcommand{\qSeqPagodaStem}{一百零八塔位于宁夏回族自治区青铜峡市，以其独特的建筑格局和深远的历史文化闻名遐迩。该塔群共有 \(108\) 座塔，依山势自上而下排成 \(12\) 行，将第 \(i\) 行中塔的座数记为 \(a_i \ (i=1, 2, \cdots, 12)\)，其中 \(a_1 = 1\)，\(a_2 = a_3 = 3\)，\(a_4 = a_5 = 5\)，且 \(a_6, a_7, \cdots, a_{12}\) 是一个首项为 \(7\)，公差为 \(2\) 的等差数列。将 \(a_1, a_2, \cdots, a_{12}\) 分为 \(6\) 组，每组 \(2\) 个数。使得每组的 \(2\) 个数之和可构成一个项数为 \(6\) 且公差为 \(d \ (d > 0)\) 的等差数列。则 \(d = \hfill （\quad）\)}

\newcommand{\qSeqPagodaOptions}{%
  \mcoptions[4]{\(2\)}{\(4\)}{\(6\)}{\(8\)}%
}

\newcommand{\qSeqPagodaAnswer}{B}

\newcommand{\qSeqPagodaSolution}{%
  【解题思路】\\
  第一步，写出全部 12 行的塔数，明确这 12 个具体数值；第二步，计算这 12 个数的总和，从而确定分组后 6 个组和的平均值；第三步，由于所有原始数据均为奇数，每组包含 2 个数，因此每组之和必为偶数。结合等差数列的对称性和奇偶性对选项进行排查，最后进行配对验证。\\
  【解析计算】\\
  根据题意，前 5 行的塔数分别为：\\
  \(a_1 = 1\)，\(a_2 = a_3 = 3\)，\(a_4 = a_5 = 5\)。\\
  后 7 行 \(a_6, \cdots, a_{12}\) 是首项为 \(7\)，公差为 \(2\) 的等差数列，故：\\
  \(a_6 = 7\)，\(a_7 = 9\)，\(a_8 = 11\)，\(a_9 = 13\)，\(a_{10} = 15\)，\(a_{11} = 17\)，\(a_{12} = 19\)。\\
  这 12 个数分别为：\(1, 3, 3, 5, 5, 7, 9, 11, 13, 15, 17, 19\)。\\
  这 12 个数的总和为：\\
  \(S = 1 + 3 + 3 + 5 + 5 + 7 + 9 + 11 + 13 + 15 + 17 + 19 = 108\)。\\
  将这 12 个数分成 6 组，每组 2 个数，设各组和构成公差为 \(d\) 的等差数列，则这 6 个组和的平均数为 \(\dfrac{108}{6} = 18\)。\\
  因为原数列中所有的数都是奇数，任意两个奇数之和必为偶数，所以这 6 个组和必须全部是偶数。\\
  由对称性，这 6 个组和可以表示为：\\
  \(18 - \dfrac{5}{2}d\)，\(18 - \dfrac{3}{2}d\)，\(18 - \dfrac{1}{2}d\)，\(18 + \dfrac{1}{2}d\)，\(18 + \dfrac{3}{2}d\)，\(18 + \dfrac{5}{2}d\)。\\
  (1) 若 \(d = 2\)，这 6 个组和为：\(13, 15, 17, 19, 21, 23\)，均为奇数，不合题意；\\
  (2) 若 \(d = 6\)，这 6 个组和为：\(3, 9, 15, 21, 27, 33\)，均为奇数，不合题意；\\
  (3) 若 \(d = 8\)，最小的组和为 \(18 - \dfrac{5}{2} \cdot 8 = -2\)，为负数，不合题意；\\
  (4) 若 \(d = 4\)，这 6 个组和为：\(8, 12, 16, 20, 24, 28\)，均为偶数。\\
  我们来验证这组和能否通过原 12 个数配对得到：\\
  \(1 + 7 = 8\)，\\
  \(3 + 9 = 12\)，\\
  \(3 + 13 = 16\)，\\
  \(5 + 15 = 20\)，\\
  \(5 + 19 = 24\)，\\
  \(11 + 17 = 28\)。\\
  恰好配对成功，各组之和构成公差为 4 的等差数列。\\
  故选 B。%
}

\newcommand{\qSeqPagodaDiagnosis}{%
  学生在处理等差数列分组求和时，容易陷入复杂的代数方程，忽略了“两个奇数之和为偶数”这一简单的数学常识。利用奇偶性排除选项可以大大提高解题速度。%
}

\newcommand{\qSeqPagodaTeachingNotes}{%
  本题以传统文化背景引入，考查等差数列的性质及逻辑推理能力。讲评时应指导学生利用整体思维：先求和，再定均值，进而利用奇偶性排查选项，最后进行快速连线配对验证。%
}
